\chapter{Implementation}
\label{ch:implementation}

\section{Project Structure}
\label{sec:project_structure}

The project follows a layered module design that mirrors the theoretical architecture described in
Chapter~\ref{ch:architecture}. Each subdirectory beneath \texttt{src/} maps to exactly one
responsibility in the data pipeline: the symbolic layer builds and maintains the knowledge graph,
the neural layer scores and distills salient changes, the agent layer routes distilled context to
specialized LLM agents, and the benchmark layer provides reproducible performance comparisons.
Keeping these concerns separate makes each component independently testable and allows the pipeline
depth to be adjusted by swapping or bypassing individual modules without touching the others.

\begin{lstlisting}[language=bash, caption={Project directory structure}]
ntcc/
  app.py                      # Streamlit entry point (legacy root launcher)
  config.py                   # API keys, thresholds, model parameters
  src/
    symbolic/                 # graph_builder.py, digital_twin.py
    neural/                   # saliency.py, distiller.py, semantic_delta.py
    agents/                   # base_agent.py, analyst_agent.py, orchestrator.py
    benchmarks/               # runner.py, baseline.py, distilled.py, ablation.py
    data_generator/           # generator.py (synthetic event stream)
    dashboard/                # app.py (four-page Streamlit interface)
  results/                    # Benchmark output JSON and PNG charts
  scripts/                    # Figure generation for the paper
  paper/
    figures/                  # PDF and PNG figures used by LaTeX
    chapters/                 # One .tex file per chapter
    main.tex                  # Top-level LaTeX document
\end{lstlisting}

The \texttt{config.py} module at the project root centralises every tunable constant: the Gemini
API key, the Gemini endpoint URL, the saliency threshold (default 0.35), and the high-saliency
threshold (default 0.65). Centralising configuration prevents the magic-number problem common in
research prototypes and makes threshold-sensitivity experiments straightforward to run without
editing multiple files.

\section{Symbolic Layer}
\label{sec:symbolic_layer}

The symbolic layer is the entry point for all raw data. Its sole responsibility is to convert an
unordered batch of business events into a structured directed graph, referred to throughout the
codebase as the \emph{digital twin}. The \texttt{build\_graph} function in
\texttt{src/symbolic/graph\_builder.py} accepts a list of event dictionaries and returns a
\texttt{networkx.DiGraph} object whose nodes represent business entities (products, suppliers,
regions, warehouses) and whose edges capture the relationships between them.

\begin{lstlisting}[language=Python, caption={graph\_builder.py: build\_graph and anomaly normalization (excerpt)}]
def build_graph(events: list[dict]) -> nx.DiGraph:
    """Build a knowledge graph from a batch of business events.

    Node types: product, supplier, region, warehouse
    Edge types: supplies, sold_in, stored_at, ships_to
    """
    G = nx.DiGraph()

    product_stats = defaultdict(lambda: {"total_sales": 0, "total_quantity": 0,
                                          "event_count": 0, "stockouts": 0,
                                          "demand_spikes": 0, "quality_issues": 0})
    supplier_stats = defaultdict(lambda: {"delays": 0, "quality_issues": 0,
                                           "shipments": 0})
    edge_weights = defaultdict(float)

    for event in events:
        product, supplier = event["product"], event["supplier"]
        region,  warehouse = event["region"],  event["warehouse"]
        etype = event["event_type"]

        for node, ntype, extras in [
            (product,   "product",   {"events": 0, "anomaly_score": 0.0}),
            (supplier,  "supplier",  {"events": 0, "reliability":   1.0}),
            (region,    "region",    {"events": 0}),
            (warehouse, "warehouse", {"events": 0, "anomaly_score": 0.0}),
        ]:
            if not G.has_node(node):
                G.add_node(node, node_type=ntype, **extras)
            G.nodes[node]["events"] += 1

        _ensure_edge(G, supplier,  product,   "supplies",   edge_weights)
        _ensure_edge(G, product,   region,    "sold_in",    edge_weights)
        _ensure_edge(G, product,   warehouse, "stored_at",  edge_weights)
        _ensure_edge(G, supplier,  warehouse, "ships_to",   edge_weights)

        if etype == "stockout":
            G.nodes[product]["anomaly_score"]   += 0.5
            G.nodes[warehouse]["anomaly_score"] += 0.3
        elif etype == "demand_spike":
            G.nodes[product]["anomaly_score"] += (
                0.3 * event.get("spike_factor", 1))

    # Sigmoid normalisation: maps raw anomaly accumulation to [0, 1]
    for node, data in G.nodes(data=True):
        if "anomaly_score" in data and data.get("events", 0) > 0:
            rate = data["anomaly_score"] / max(data["events"], 1)
            data["anomaly_score"] = round(
                min(1.0, max(0.0, 2.0 / (1.0 + math.exp(-5.0 * rate)) - 1.0)), 4)

    _propagate_risk(G)
    return G
\end{lstlisting}

Several design decisions embedded in this function are worth noting. First, every node receives a
\texttt{node\_type} attribute at creation time; downstream modules use this attribute to apply
type-specific importance weights without performing any additional graph traversal. Second,
anomaly scores are accumulated as raw floats during event processing and then normalised via a
sigmoid-like mapping before the function returns. This means the score for a product that
experienced ten stockouts is not ten times the score for a product that experienced one; instead,
the mapping applies diminishing returns so that a single severe event already pushes the score
toward one. Third, the final call to \texttt{\_propagate\_risk} ensures that supplier
reliability degradation cascades downstream to the products that supplier feeds, capturing the
second-order effects that naive event counting would miss entirely.

\section{Neural Layer}
\label{sec:neural_layer}

Once two consecutive graph snapshots are available, the neural layer computes the \emph{semantic
delta}---the set of structural changes between the two graphs---and then scores each change by
its business saliency. The scoring logic lives in \texttt{src/neural/saliency.py}. Rather than
applying a flat threshold to every change type, the scorer combines three independent signals: a
node-type weight (suppliers outrank products, which outrank warehouses and regions), a centrality
boost derived from the current graph topology, and an attribute-level change magnitude. Only
changes whose combined score exceeds \texttt{SALIENCY\_THRESHOLD} are forwarded to the
distiller; the rest are silently discarded, forming the noise-filtering layer described in
Section~\ref{sec:methodology_distillation}.

\begin{lstlisting}[language=Python, caption={saliency.py: score\_delta with centrality boost (excerpt)}]
NODE_TYPE_WEIGHTS = {
    "supplier":  1.2,
    "product":   1.0,
    "warehouse": 0.8,
    "region":    0.6,
}

def score_delta(delta: dict, g_curr: nx.DiGraph = None) -> list[dict]:
    """Score each element of the semantic delta by structural saliency."""
    scored   = []
    centrality = {}

    if g_curr is not None and g_curr.number_of_nodes() > 0:
        try:
            degree_cent  = nx.degree_centrality(g_curr)
            betweenness  = nx.betweenness_centrality(g_curr)
            for node in g_curr.nodes():
                # Combined centrality scaled to [0, 1]
                centrality[node] = min(
                    1.0,
                    (degree_cent.get(node, 0) + betweenness.get(node, 0)) / 2 * 3
                )
        except Exception:
            pass

    for item in delta["changed_nodes"]:
        node        = item["node"]
        node_weight = NODE_TYPE_WEIGHTS.get(item["node_type"], 0.5)
        cent_boost  = 1.0 + centrality.get(node, 0) * 0.5  # up to 1.5x

        max_change_score = max(
            (_attribute_saliency(attr, vals)
             for attr, vals in item["changes"].items()),
            default=0
        )
        score = min(1.0, max_change_score * node_weight * cent_boost)

        scored.append({
            "type":       "changed_node",
            "entity":     node,
            "node_type":  item["node_type"],
            "changes":    item["changes"],
            "saliency":   round(score, 3),
            "centrality": round(centrality.get(node, 0), 3),
        })

    # Broken edges are more critical than new ones
    for item in delta["removed_edges"]:
        src, tgt = item["source"], item["target"]
        score = 0.4 + (centrality.get(src, 0.3) + centrality.get(tgt, 0.3)) * 0.3
        scored.append({
            "type":     "removed_edge",
            "entity":   src,
            "saliency": round(min(1.0, score), 3),
            "reason":   f"Broken relationship: {src} -> {tgt}",
        })

    return scored
\end{lstlisting}

The centrality boost is the most consequential design choice in this module. A product node that
sits at the intersection of many supplier-warehouse paths receives a higher weight than an
isolated peripheral node showing the same absolute anomaly score. This mirrors how a supply chain
analyst would naturally prioritise: a disruption at a hub product affects far more downstream
operations than an equivalent disruption at a niche SKU. The betweenness centrality component
captures this bridging role explicitly, and the degree centrality component captures sheer
connectivity. Both are averaged and scaled to a boost factor that ranges from 1.0 (no topology
contribution) to 1.5 (maximally central node), keeping the final score bounded to $[0, 1]$.

\section{Agent Layer}
\label{sec:agent_layer}

The agent layer wraps Google Gemini API calls in a clean class hierarchy. \texttt{BaseAgent} in
\texttt{src/agents/base\_agent.py} handles HTTP communication, exponential backoff on 429 and
5xx errors, and token-usage accounting. Concrete agent subclasses inherit from \texttt{BaseAgent}
and override only the system prompt and the task-specific method that formats the user turn.
Three specialised agents are implemented: \texttt{AnalystAgent} for business risk and
opportunity analysis, \texttt{RiskAgent} for quantified supply chain risk assessment, and
\texttt{ForecastAgent} for short-term trend prediction.

\begin{lstlisting}[language=Python, caption={analyst\_agent.py: AnalystAgent role definition and inference (excerpt)}]
class AnalystAgent(BaseAgent):
    """Business analyst that interprets distilled BI briefings."""

    def __init__(self):
        super().__init__(
            name="Business Analyst",
            role="analyst",
            system_prompt=(
                "You are a senior Business Intelligence analyst. You receive distilled "
                "briefings about changes in the business knowledge graph. Your job is to:\n"
                "1. Identify the most critical business risks\n"
                "2. Spot opportunities hidden in the data\n"
                "3. Recommend specific actions\n"
                "4. Assess supply chain health\n\n"
                "Be concise and actionable. Use bullet points. "
                "Prioritize by business impact."
            )
        )

    def analyze(self, briefing: str) -> dict:
        """Analyze a distilled briefing and produce insights."""
        prompt = (
            f"Analyze this business intelligence briefing and provide:\n"
            f"1. TOP RISKS (ranked by severity)\n"
            f"2. OPPORTUNITIES\n"
            f"3. RECOMMENDED ACTIONS\n"
            f"4. SUPPLY CHAIN HEALTH SCORE (1-10)\n\n"
            f"{briefing}"
        )
        return self.query(prompt, temperature=0.4)


class RiskAgent(BaseAgent):
    """Specialized risk assessment agent."""

    def assess_risk(self, briefing: str) -> dict:
        prompt = (
            f"Assess the risks in this briefing. For each risk:\n"
            f"- Likelihood (Low/Medium/High)\n"
            f"- Impact (Low/Medium/High/Critical)\n"
            f"- Recommended mitigation\n\n"
            f"{briefing}"
        )
        return self.query(prompt, temperature=0.3)
\end{lstlisting}

The temperature values are set deliberately low (0.3--0.4) for the analyst and risk agents because
reproducibility and factual grounding matter more than creative variation in a business
intelligence context. The forecast agent uses a slightly higher temperature of 0.5 because
plausible trend extrapolation benefits from a modest degree of generative diversity. All three
agents receive the same distilled briefing text as their user turn, which means their outputs can
be directly compared and, where appropriate, merged by the orchestrator into a single composite
report. The orchestrator in \texttt{src/agents/orchestrator.py} coordinates this fan-out and
fan-in pattern, invoking each agent in sequence and aggregating the results.

\section{Streamlit Dashboard}
\label{sec:dashboard}

The Streamlit dashboard in \texttt{src/dashboard/app.py} provides an interactive four-page
interface for exploring every stage of the pipeline. The \emph{Home} page presents the full
pipeline as a live demonstration: the user selects an event count and anomaly ratio, the system
generates a synthetic event batch, runs the complete distillation pipeline, and displays the
resulting Business Intelligence Briefing alongside compression and latency metrics. The
\emph{How It Works} page renders an annotated walkthrough of the data flow from raw events to
distilled text. The \emph{Knowledge Graph} page renders the current digital twin as an
interactive Plotly network diagram. The \emph{Comparison} page places the naive (baseline) and
distilled approaches side by side, including the full raw context that the baseline would have
sent to the LLM and the condensed briefing that the distilled approach sends instead.

\begin{figure}[H]
  \centering
  
\includegraphics[width=0.9\textwidth]{review_graph}
  \caption{Knowledge graph visualization in the Streamlit dashboard. Nodes are colored by
           type (product, supplier, warehouse, region) and sized by degree centrality.}
  \label{fig:review_graph}
\end{figure}

The knowledge graph page shown in Figure~\ref{fig:review_graph} renders the \texttt{NetworkX}
graph using Plotly's scatter-plus-line trace approach. Node positions are computed with the
Kamada--Kawai layout algorithm, which tends to place well-connected hub nodes near the center
and peripheral nodes toward the edges of the canvas---a spatial arrangement that visually
reinforces the centrality-based scoring decisions made in the neural layer. Nodes are colored by
type using a four-color scheme (teal for products, orange for suppliers, blue for warehouses,
green for regions) and sized proportionally to their degree centrality, so high-risk hubs are
immediately legible at a glance without requiring the user to inspect individual node metadata.

\begin{figure}[H]
  \centering
  \includegraphics[width=0.9\textwidth]{review_livedemo_result}
  \caption{Live pipeline demo page showing the distilled Business Intelligence Briefing
           generated from a synthetic event batch.}
  \label{fig:livedemo}
\end{figure}

Figure~\ref{fig:livedemo} shows the result panel of the live demo page after a complete pipeline
run. The panel displays the distilled briefing text alongside four key metrics: total input
tokens sent to the LLM, compression ratio relative to the naive baseline, end-to-end latency in
seconds, and the share of raw events filtered out as noise. In a representative run with 500
events at a 20\% anomaly ratio, the pipeline compressed the input from approximately 12\,400
tokens (baseline) to 680 tokens (distilled), a compression ratio of $18.2\times$, while
retaining all anomalous events in the briefing output.

\section{Benchmark Framework}
\label{sec:benchmark}

Reproducible performance measurement requires a framework that runs multiple trials per
configuration, averages results across runs to reduce variance, and serialises outputs to JSON
for later analysis and figure generation. The benchmark runner in
\texttt{src/benchmarks/runner.py} implements this framework. It sweeps across a configurable
list of event counts (default: 100, 300, 500, 1000 events), executes three independent
trials at each size, and computes per-trial comparisons as well as per-size averages. The
\texttt{run\_baseline} function in \texttt{src/benchmarks/baseline.py} sends the full raw event
list directly to the LLM; the \texttt{run\_distilled} function in
\texttt{src/benchmarks/distilled.py} runs the complete neuro-symbolic pipeline and sends only
the resulting briefing.

\begin{lstlisting}[language=Python, caption={runner.py: multi-trial benchmark loop (excerpt)}]
def run_benchmark(
    event_counts:  list[int] = [100, 300, 500, 1000],
    anomaly_ratio: float     = 0.2,
    runs_per_size: int       = 3,
) -> dict:
    """Run benchmarks across different data sizes with multiple runs."""
    os.makedirs(RESULTS_DIR, exist_ok=True)
    results = {"baseline": [], "distilled": [], "comparisons": [], "averaged": []}

    for count in event_counts:
        size_comparisons = []

        for run in range(runs_per_size):
            events_t0 = generate_batch(count, anomaly_ratio=0.05)
            events_t1 = generate_batch(count, anomaly_ratio=anomaly_ratio)

            baseline_result  = run_baseline(events_t1)
            distilled_result = run_distilled(events_t0, events_t1)

            comparison = {
                "event_count":            count,
                "run":                    run,
                "baseline_input_tokens":  baseline_result["input_tokens"],
                "distilled_input_tokens": distilled_result["input_tokens"],
                "compression_ratio":      distilled_result["compression_ratio"],
                "baseline_latency":       baseline_result["latency_seconds"],
                "distilled_latency":      distilled_result["latency_seconds"],
                "pipeline_overhead":      distilled_result["pipeline_processing_seconds"],
                "noise_filtered_pct":     round(
                    distilled_result["reduction_ratio"] * 100, 1),
            }
            results["comparisons"].append(comparison)
            size_comparisons.append(comparison)
            time.sleep(1)

        avg = _average_comparisons(size_comparisons, count)
        results["averaged"].append(avg)

    results_path = os.path.join(RESULTS_DIR, "benchmark_results.json")
    with open(results_path, "w") as f:
        json.dump(results, f, indent=2, default=str)

    generate_charts(results)
    return results
\end{lstlisting}

The one-second \texttt{time.sleep} between calls is a deliberate rate-limit courtesy to the
Gemini API; it does not affect wall-clock benchmark fairness because both the baseline and
distilled calls are subject to the same pause. The call to \texttt{generate\_charts} at the end
produces the bar charts and compression-ratio plots included in Chapter~\ref{ch:results}. The
JSON output is also consumed by \texttt{scripts/generate\_report\_figures.py}, which regenerates
all publication-quality PDF figures from the raw numeric results, making it straightforward to
update figures whenever additional benchmark runs are completed.
